\documentclass{article}
\usepackage{amsmath}
\usepackage[top=3cm, bottom=3cm, outer=3cm, inner=3cm]{geometry}
\usepackage{multicol}
\usepackage{graphicx}
\usepackage{booktabs}
\usepackage{pdflscape}
\usepackage{url}
\usepackage[natbibapa]{apacite}
\usepackage{array}
\newcolumntype{x}[1]{>{\centering\arraybackslash\hspace{0pt}}p{#1}}
\usepackage{hyperref}
\usepackage{pdfpages}
\usepackage{multirow}
\usepackage[normalem]{ulem}
\useunder{\uline}{\ul}{}
\usepackage{svg}
\usepackage[T1]{fontenc}
\usepackage{xcolor}
\usepackage{listings}
\lstdefinestyle{ascii-tree}{
    literate={├}{|}1 {─}{--}1 {└}{+}1 
  }
\lstset{basicstyle=\ttfamily,
  showstringspaces=false,
  commentstyle=\color{red},
  keywordstyle=\color{blue}
}
\usepackage{caption}
\usepackage{subcaption}
\usepackage{float}
\newcolumntype{M}[1]{>{\centering\arraybackslash}m{#1}}
\newcolumntype{N}{@{}m{0pt}@{}}


%%%%%%%%%%%%%%%%%%%%%%%%%%%%%%%%%%%%%%%%%%%%%%%%%%%%%%%%%%%%%%%%%%%%%%%%%%%%
%%%%%%%%%%%%%%%%%%%%%%%%%%%%%%%%%%%%%%%%%%%%%%%%%%%%%%%%%%%%%%%%%%%%%%%%%%%%
\newcommand{\itemEmail}{ccapatinta@unsa.edu.pe}
\newcommand{\itemStudent}{Camilo B.A. Capatinta Almanza}
\newcommand{\itemCourse}{Análisis y Diseño de Algoritmos}
\newcommand{\itemCourseCode}{1702231}
\newcommand{\itemSemester}{IV}
\newcommand{\itemUniversity}{Universidad Nacional de San Agustín de Arequipa}
\newcommand{\itemFaculty}{Facultad de Ingeniería de Producción y Servicios}
\newcommand{\itemDepartment}{Departamento Académico de Ingeniería de Sistemas e Informática}
\newcommand{\itemSchool}{Escuela Profesional de Ingeniería de Sistemas}
\newcommand{\itemAcademic}{2026 - B}
\newcommand{\itemInput}{22/10/2026}
\newcommand{\itemOutput}{28/10/2026}
\newcommand{\itemDocente}{Harvey A. Norman Patrick}
\newcommand{\itemTheme}{Algoritmos recursivos}
\newcommand{\itemPracticeNumber}{04}
%%%%%%%%%%%%%%%%%%%%%%%%%%%%%%%%%%%%%%%%%%%%%%%%%%%%%%%%%%%%%%%%%%%%%%%%%%%%
%%%%%%%%%%%%%%%%%%%%%%%%%%%%%%%%%%%%%%%%%%%%%%%%%%%%%%%%%%%%%%%%%%%%%%%%%%%%

\usepackage[utf8]{inputenc}
\usepackage[english,spanish]{babel}
\AtBeginDocument{\selectlanguage{spanish}}
\renewcommand{\figurename}{Figura}
\renewcommand{\refname}{Referencias}
\AtBeginDocument{%
	\renewcommand\tablename{Tabla}
}

\usepackage{fancyhdr}
\pagestyle{fancy}
\fancyhf{}
\setlength{\headheight}{40.51407pt}
\renewcommand{\headrulewidth}{1pt}
\renewcommand{\footrulewidth}{1pt}
\fancyhead[L]{\raisebox{-0.2\height}{\includegraphics[width=3cm]{img/logo_episunsa.png}}}
\fancyhead[C]{\fontsize{7}{7}\selectfont \itemUniversity \\ \itemFaculty \\ \itemDepartment \\ \itemSchool \\ \textbf{\itemCourse}}
\fancyhead[R]{\raisebox{-0.2\height}{\includegraphics[width=1.2cm]{img/logo_abet.png}}}
\fancyfoot[L]{Estudiante Capatinta Camilo}
\fancyfoot[C]{\itemCourse}
\fancyfoot[R]{Página \thepage}

\newcommand{\landscapeheader}{%
    \begin{tabular*}{\linewidth}{@{\extracolsep{\fill}}c c c@{}}
        \includegraphics[width=3cm]{img/logo_episunsa.png} &
        {\fontsize{7}{8}\selectfont\shortstack[c]{
        \itemUniversity\\
        \itemFaculty\\
        \itemDepartment\\
        \itemSchool\\
        \textbf{\itemCourse}}} &
        \includegraphics[width=1.2cm]{img/logo_abet.png}
    \end{tabular*}
    \par\vspace{0.15cm}\hrule\vspace{0.4cm}
}

\newcommand{\landscapefooter}{%
    \vfill
    \noindent\hrule\vspace{0.15cm}
    \noindent Estudiante Capatinta Camilo\hfill\itemCourse\hfill Página \thepage
}

% Para el código fuente
\usepackage{colortbl}
\definecolor{dkgreen}{rgb}{0,0.6,0}
\definecolor{gray}{rgb}{0.5,0.5,0.5}
\definecolor{mauve}{rgb}{0.58,0,0.82}
\definecolor{codebackground}{rgb}{0.95, 0.95, 0.92}
\definecolor{tablebackground}{rgb}{0.8, 0, 0}

\lstset{frame=tb,
	language=bash,
	aboveskip=3mm,
	belowskip=3mm,
	showstringspaces=false,
	columns=flexible,
	basicstyle={\small\ttfamily},
	numbers=none,
	numberstyle=\tiny\color{gray},
	keywordstyle=\color{blue},
	commentstyle=\color{dkgreen},
	stringstyle=\color{mauve},
	breaklines=true,
	breakatwhitespace=true,
	tabsize=3,
	backgroundcolor= \color{codebackground},
}

% Título del índice en español y enlaces en el índice
\addto\captionsspanish{\renewcommand{\contentsname}{Índice}}
\hypersetup{
    colorlinks=true,
    allcolors=black,
    linktoc=all        % toda la línea del índice es clicable
}

% Evita que natbib imprima su propio título "Referencias"
% (usamos nuestra propia \section para que aparezca numerada en el índice)
\makeatletter
\renewcommand{\bibsection}{}
\makeatother

\begin{document}
	
	\vspace*{10px}
	
	\begin{center}	
		\fontsize{17}{17} \textbf{ Actividad \itemPracticeNumber}
	\end{center}
	\centerline{\textbf{\Large Tema: \itemTheme}}
	%\vspace*{0.5cm}	

	\begin{flushright}
		\begin{tabular}{|M{2.5cm}|N|}
			\hline 
			\rowcolor{tablebackground}
			\color{white} \textbf{Nota}  \\
			\hline 
			     \\[30pt]
			\hline 			
		\end{tabular}
	\end{flushright}	

	\begin{table}[h!]
		\begin{tabular}{|x{4.7cm}|x{4.8cm}|x{4.8cm}|}        
			\hline 
			\rowcolor{tablebackground}
			\color{white} \textbf{Estudiante} & \color{white}\textbf{Escuela}  & \color{white}\textbf{Asignatura}   \\
			\hline 
			{\itemStudent \par \itemEmail} & \itemSchool & {\itemCourse \par Semestre: \itemSemester \par Código: \itemCourseCode}     \\
			\hline 			
		\end{tabular}
	\end{table}		
	
	\begin{table}[H]
		\begin{tabular}{|x{4.7cm}|x{4.8cm}|x{4.8cm}|}
			\hline 
			\rowcolor{tablebackground}
			\color{white}\textbf{Docente} & \color{white}\textbf{Tema}  & \color{white}\textbf{Duración}   \\
			\hline 
			\itemDocente & \itemTheme & 04 horas   \\
			\hline 
		\end{tabular}
	\end{table}
	
	\begin{table}[H]
		\begin{tabular}{|x{4.7cm}|x{4.8cm}|x{4.8cm}|}
			\hline 
			\rowcolor{tablebackground}
			\color{white}\textbf{Semestre académico} & \color{white}\textbf{Fecha de inicio}  & \color{white}\textbf{Fecha de entrega}   \\
			\hline 
			\itemAcademic & \itemInput &  \itemOutput  \\
			\hline 
		\end{tabular}
	\end{table}
    \clearpage
    \tableofcontents
    \clearpage

    \section{Actividades}
 
\subsection{Objetivo}
 
Investigar el diseño de algoritmos recursivos: su definición, los elementos
que debe tener todo algoritmo recursivo para ser correcto, los tipos de
recursión, el proceso para diseñar uno a partir de un problema, su relación
con otras técnicas de diseño (divide y vencerás, backtracking, programación
dinámica) y sus ventajas y desventajas frente a una solución iterativa.
 
% ------------------------------------------------------------
\subsection{Definición}
 
Un \textit{algoritmo recursivo} es aquel que resuelve un problema invocándose
a sí mismo sobre una o más instancias más pequeñas del mismo problema, hasta
llegar a una instancia lo bastante simple como para resolverse directamente
\citep{clrs}. La recursión traslada al algoritmo la misma idea que una
definición recursiva en matemáticas: un objeto se define en términos de una
versión más simple de sí mismo \citep{wirth1976}.
 
% ------------------------------------------------------------
\subsection{Elementos de un algoritmo recursivo correcto}
 
Para que la recursión termine y produzca el resultado correcto, todo
algoritmo recursivo necesita los siguientes elementos \citep{levitin2012}:
 
\begin{itemize}
  \item \textbf{Caso base (o caso trivial):} una o más instancias del
        problema que se resuelven directamente, sin nuevas llamadas
        recursivas. Es la condición de parada de la recursión.
  \item \textbf{Caso recursivo:} la regla que expresa la solución del
        problema en términos de la solución de una o más instancias más
        pequeñas del mismo problema.
  \item \textbf{Progreso hacia el caso base:} cada llamada recursiva debe
        operar sobre una instancia estrictamente más cercana al caso base
        (por ejemplo, un tamaño de entrada menor), de modo que el proceso
        termine en un número finito de pasos.
  \item \textbf{Combinación de resultados (cuando aplica):} cuando el
        problema se divide en más de una subinstancia, el algoritmo debe
        combinar las soluciones parciales para construir la solución del
        problema original.
\end{itemize}
 
Si falta el caso base, o si las llamadas recursivas no se acercan a él, el
algoritmo entra en una recursión infinita que termina por agotar la memoria
disponible para las llamadas \citep{brassardbratley1996}.
 
% ------------------------------------------------------------
\subsection{El principio de inducción y la correctitud de la recursión}
 
La correctitud de un algoritmo recursivo se demuestra de la misma forma que
una demostración por \textit{inducción matemática} \citep{clrs}:
 
\begin{itemize}
  \item Se demuestra que el algoritmo es correcto para el caso base (caso
        base de la inducción).
  \item Se asume que el algoritmo es correcto para instancias más pequeñas
        que la actual (hipótesis inductiva).
  \item Se demuestra que, asumiendo lo anterior, el algoritmo también es
        correcto para la instancia actual (paso inductivo).
\end{itemize}
 
Esta correspondencia explica por qué la recursión resulta natural para
problemas que ya tienen una definición inductiva, como el factorial, los
números de Fibonacci o el recorrido de estructuras definidas recursivamente
(listas, árboles).
 
% ------------------------------------------------------------
\subsection{Proceso de diseño de un algoritmo recursivo}
 
Diseñar un algoritmo recursivo sigue, en general, los siguientes pasos
\citep{levitin2012}:
 
\begin{itemize}
  \item \textbf{Paso 1:} identificar uno o más casos base, es decir, las
        instancias más pequeñas o simples del problema que se pueden
        resolver sin recursión.
  \item \textbf{Paso 2:} encontrar una forma de reducir el problema original
        a una o más instancias más pequeñas del mismo problema (la relación
        de recurrencia del algoritmo).
  \item \textbf{Paso 3:} expresar la solución del problema en función de la
        solución de esas instancias más pequeñas, incluyendo cómo combinar
        los resultados si hay más de una llamada recursiva.
  \item \textbf{Paso 4:} verificar que cada llamada recursiva progresa hacia
        el caso base (por ejemplo, que el tamaño de la entrada disminuye en
        cada llamada).
  \item \textbf{Paso 5:} analizar el algoritmo resultante planteando su
        relación de recurrencia de tiempo de ejecución y resolviéndola (por
        ejemplo, con el Teorema Maestro cuando aplica).
\end{itemize}
 
% ------------------------------------------------------------
\subsection{Tipos de recursión}
 
\begin{itemize}
  \item \textbf{Recursión directa:} un procedimiento se llama a sí mismo de
        forma explícita, como en el cálculo del factorial.
  \item \textbf{Recursión indirecta (o mutua):} un procedimiento $A$ llama a
        un procedimiento $B$, y $B$ llama de nuevo a $A$, formando un ciclo
        de llamadas entre dos o más procedimientos.
  \item \textbf{Recursión lineal:} cada llamada genera como máximo una nueva
        llamada recursiva (por ejemplo, el factorial o la búsqueda binaria).
        El árbol de llamadas resultante es, de hecho, una lista.
  \item \textbf{Recursión múltiple (o en árbol):} cada llamada genera más de
        una llamada recursiva, como en el cálculo ingenuo de Fibonacci o en
        el recorrido de un árbol binario. El árbol de llamadas crece en
        forma de árbol, no de lista.
  \item \textbf{Recursión anidada:} uno de los argumentos de la llamada
        recursiva es, a su vez, el resultado de otra llamada recursiva (por
        ejemplo, la función de Ackermann).
  \item \textbf{Recursión de cola (\emph{tail recursion}):} la llamada
        recursiva es la última operación que realiza el procedimiento, sin
        trabajo pendiente después de que esta retorne. Muchos compiladores e
        intérpretes pueden optimizar este caso convirtiéndolo internamente
        en un bucle, evitando el crecimiento de la pila de llamadas
            \citep{brassardbratley1996}.
\end{itemize}
 
% ------------------------------------------------------------
\subsection{La pila de llamadas y el costo en memoria}
 
Cada llamada recursiva que todavía no ha retornado ocupa un marco de
activación (\emph{activation record}) en la pila de llamadas del programa,
con sus parámetros, variables locales y la dirección de retorno
\citep{clrs}. Esto tiene dos consecuencias prácticas:
 
\begin{itemize}
  \item La \textbf{profundidad máxima de la recursión} (el número de
        llamadas pendientes en el peor caso) determina la complejidad
        espacial adicional que introduce la recursión, más allá del espacio
        que usaría una solución iterativa equivalente.
  \item Una recursión sin caso base, o cuyo progreso hacia el caso base es
        incorrecto, agota la pila de llamadas y produce un error de
        \emph{stack overflow} (desbordamiento de pila).
\end{itemize}
 
% ------------------------------------------------------------
\subsection{Recursión frente a iteración}
 
Todo algoritmo recursivo puede reescribirse como un algoritmo iterativo
equivalente (y viceversa), ya que ambos modelos de cómputo tienen el mismo
poder expresivo \citep{brassardbratley1996}. La elección entre uno u otro
es una decisión de diseño:
 
\begin{itemize}
  \item \textbf{A favor de la recursión:} el código suele ser más corto, más
        cercano a la definición matemática o a la estructura de datos del
        problema (listas, árboles, grafos), y más fácil de demostrar
        correcto mediante inducción.
  \item \textbf{A favor de la iteración:} evita el costo en memoria y en
        tiempo de las llamadas de función (creación y destrucción de marcos
        de activación), y evita el riesgo de \emph{stack overflow} en
        entradas grandes cuando el lenguaje o el entorno no garantiza
        optimización de recursión de cola.
  \item En la práctica, se prefiere la recursión cuando el problema tiene una
        estructura naturalmente recursiva y la profundidad de recursión es
        manejable, y la iteración cuando el rendimiento o el límite de pila
        son críticos.
\end{itemize}
 
% ------------------------------------------------------------
\subsection{Relación con otras técnicas de diseño de algoritmos}
 
La recursión es el mecanismo que hace posible varias técnicas de diseño de
algoritmos más generales \citep{clrs}:
 
\begin{itemize}
  \item \textbf{Divide y vencerás:} divide el problema en subproblemas
        independientes del mismo tipo, los resuelve recursivamente y combina
        sus soluciones (por ejemplo, Merge Sort, Quicksort, la multiplicación
        de Karatsuba).
  \item \textbf{Backtracking (vuelta atrás):} explora recursivamente el
        espacio de soluciones de un problema, retrocediendo cuando una
        decisión parcial no puede llevar a una solución válida (por ejemplo,
        el problema de las $N$ reinas o la resolución de laberintos).
  \item \textbf{Programación dinámica (enfoque top-down con memoización):}
        parte de la misma relación de recurrencia que una solución recursiva
        directa, pero almacena los resultados de las subinstancias ya
        resueltas para no recalcularlas, lo que puede reducir drásticamente
        el tiempo de ejecución de una recursión múltiple con subproblemas
        repetidos.
\end{itemize}
 
% ------------------------------------------------------------
\subsection{Ejemplos de diseño recursivo}
 
\subsubsection{Ejemplo 1: Factorial (recursión lineal)}
 
\textbf{Caso base:} $0! = 1$.

\textbf{Caso recursivo:} $n! = n \cdot (n-1)!$ para $n \geq 1$.
 
Cada llamada reduce el problema en una unidad ($n \to n-1$), por lo que el
progreso hacia el caso base es inmediato. Es un ejemplo de recursión lineal:
el árbol de llamadas es, en realidad, una cadena de $n$ llamadas.
 
\subsubsection{Ejemplo 2: Sucesión de Fibonacci (recursión múltiple)}
 
\textbf{Caso base:} $F(0) = 0$, $F(1) = 1$.

\textbf{Caso recursivo:} $F(n) = F(n-1) + F(n-2)$ para $n \geq 2$.
 
Cada llamada genera dos llamadas recursivas, por lo que el árbol de llamadas
crece de forma exponencial y muchas subinstancias se recalculan varias
veces (por ejemplo, $F(2)$ se calcula repetidamente al calcular $F(5)$). Este
es el ejemplo clásico que motiva la programación dinámica con memoización
como mejora directa sobre el diseño recursivo ingenuo.
 
\subsubsection{Ejemplo 3: Torres de Hanoi (divide y vencerás recursivo)}
 
\textbf{Caso base:} mover $0$ discos no requiere ninguna acción.

\textbf{Caso recursivo:} para mover $n$ discos del poste origen al poste
destino usando un poste auxiliar: mover $n-1$ discos del origen al auxiliar,
mover el disco restante del origen al destino, y mover los $n-1$ discos del
auxiliar al destino.
 
Aquí la reducción no es solo en el tamaño del problema ($n \to n-1$), sino
que además se necesitan dos llamadas recursivas para resolver una sola
instancia, ilustrando cómo la recursión expresa de forma directa una
estrategia de divide y vencerás.
\subsubsection{Ejemplo 4: Búsqueda binaria (recursión de cola)}
 
\textbf{Caso base:} el subarreglo está vacío (no se encontró el elemento), o
el elemento central es igual al buscado.

\textbf{Caso recursivo:} si el elemento buscado es menor que el elemento
central, buscar recursivamente en la mitad izquierda; si es mayor, buscar en
la mitad derecha.
 
La llamada recursiva es la última operación del algoritmo y no queda trabajo
pendiente después de que retorne, por lo que es un ejemplo de recursión de
cola: puede convertirse directamente en un bucle \texttt{while} sin cambiar
su comportamiento.
 
% ------------------------------------------------------------
\subsection{Resumen}
 
\begin{center}
\small
\renewcommand{\arraystretch}{1.5}
\begin{tabular}{|l|l|l|}
\hline
\textbf{Ejemplo} & \textbf{Tipo de recursión} & \textbf{Relación de recurrencia} \\
\hline
Factorial & Lineal & $T(n) = T(n-1) + \Theta(1)$ \\
\hline
Fibonacci (ingenuo) & Múltiple (en árbol) & $T(n) = T(n-1) + T(n-2) + \Theta(1)$ \\
\hline
Torres de Hanoi & Múltiple (divide y vencerás) & $T(n) = 2T(n-1) + \Theta(1)$ \\
\hline
Búsqueda binaria & De cola & $T(n) = T(n/2) + \Theta(1)$ \\
\hline
\end{tabular}
\end{center}
 
El diseño recursivo, entonces, no es solo una forma distinta de escribir un
algoritmo: define directamente la relación de recurrencia que describe su
costo, la cual puede resolverse con los métodos revisados anteriormente
(iteración, sustitución, árbol de recursión, ecuación característica o el
Teorema Maestro) para conocer su complejidad.

\section{Referencias bibliográficas}
    \bibliographystyle{apacite}
    \nocite{*}
    \bibliography{referencias}

\end{document}